\section{Conclusion}

\label{sec:conclusion}
% ============================================================================

The HANZO token is designed to solve the fundamental tension in AI infrastructure
tokenomics: creating genuine utility demand while avoiding the velocity problem
that plagues payment-only tokens. By combining five utility dimensions---compute
payment, node staking, governance, fee reduction, and revenue sharing---with a
fixed supply and 25\% fee burn, HANZO creates multiple concurrent demand sinks
that reduce circulating supply as network usage grows.

The key properties of the model are:

\begin{enumerate}[leftmargin=1.4em]
    \item \textbf{Organic demand}: Every inference request, training job, and API
    call on Hanzo Platform requires HANZO, creating demand proportional to actual
    usage rather than speculation.
    \item \textbf{Supply reduction}: Staking, fee-reduction lock-ups, and
    permanent burns continuously reduce effective circulating supply.
    \item \textbf{Alignment}: Node operators, model providers, application
    builders, and end users all benefit from holding HANZO, aligning incentives
    across the ecosystem.
    \item \textbf{Stability}: The burn-and-stake equilibrium provides a natural
    price floor that rises with cumulative platform usage.
    \item \textbf{Governance}: DAO-controlled parameters ensure the economic model
    can adapt to changing market conditions without hard forks.
\end{enumerate}

Hanzo Network, backed by the Hanzo Platform's 14-month production track record
and the Lux Network's multi-consensus infrastructure, provides the foundation for
a sustainable AI compute economy where the HANZO token captures value
proportional to the real economic activity it enables.

% ============================================================================
% REFERENCES
% ============================================================================
\begin{thebibliography}{99}

\bibitem{hanzoplatform}
M.~Chen, D.~Wei, and Z.~Kelling.
\newblock Hanzo Platform: {PaaS} for {AI}-native application deployment.
\newblock \emph{Hanzo AI Research}, December 2021.

\bibitem{hanzoaci}
Z.~Kelling, D.~Wei, and M.~Chen.
\newblock Hanzo {ACI}: {AI} Chain Infrastructure for decentralized compute markets.
\newblock \emph{Hanzo AI Research}, December 2024.

\bibitem{hanzonetwork}
Z.~Kelling.
\newblock Hanzo Network: A {Hamiltonian} Market Maker {Layer-1} for decentralized {AI} compute and semantic learning.
\newblock \emph{Hanzo AI Research}, October 2025.

\bibitem{hanzodso}
{Hanzo AI Research}.
\newblock Decentralized Semantic Optimization: {Byzantine}-robust prior aggregation for collective model adaptation.
\newblock \emph{Hanzo AI Research}, December 2023.

\bibitem{hanzoaso}
{Hanzo AI Research}.
\newblock Active Semantic Optimization: Training-free adaptation via {Bayesian} Product-of-Experts decoding.
\newblock \emph{Hanzo AI Research}, September 2025.

\bibitem{hanzollmgateway}
{Hanzo AI Research}.
\newblock Hanzo {LLM} Gateway: Unified model serving infrastructure.
\newblock \emph{Hanzo AI Research}, 2024.

\bibitem{luxwhitepaper}
{Lux Partners}.
\newblock Lux: A scalable multi-consensus blockchain with post-quantum cryptography.
\newblock \emph{Lux Network Technical Report}, 2024.

\bibitem{mckinseyai2024}
{McKinsey Global Institute}.
\newblock The economic potential of generative {AI}: The next productivity frontier.
\newblock \emph{McKinsey Report}, June 2023.

\bibitem{buterin2014}
V.~Buterin.
\newblock A next-generation smart contract and decentralized application platform.
\newblock \emph{Ethereum White Paper}, 2014.

\bibitem{nakamoto2008}
S.~Nakamoto.
\newblock Bitcoin: A peer-to-peer electronic cash system.
\newblock 2008.

\bibitem{samani2017}
K.~Samani.
\newblock Understanding token velocity.
\newblock \emph{Multicoin Capital}, September 2017.

\bibitem{vitalikburn}
V.~Buterin, E.~Conner, R.~Dudley, M.~Slipper, and I.~Norden.
\newblock {EIP-1559}: Fee market change for {ETH 1.0} chain.
\newblock \emph{Ethereum Improvement Proposals}, April 2019.

\bibitem{fetchai2023}
{Fetch.ai}.
\newblock Fetch.ai: Autonomous economic agents.
\newblock \emph{Fetch.ai Documentation}, 2023.

\bibitem{rendernetwork}
{Render Network}.
\newblock The Render Network: Distributed {GPU} rendering.
\newblock \emph{Render Network White Paper}, 2022.

\bibitem{bittensor2023}
{Opentensor Foundation}.
\newblock Bittensor: A peer-to-peer intelligence market.
\newblock \emph{Bittensor White Paper}, 2023.

\bibitem{akash2020}
G.~Akash.
\newblock Akash Network: The decentralized cloud.
\newblock \emph{Akash Network White Paper}, 2020.

\end{thebibliography}

